% Cuerpo científico recuperado y distribuido en su residencia terminal.
\section{Clôture cofinale du critère de Riemann--Weil}

La réalisation généalogique produit une famille cofinale et naturelle de formes
positives dont la publication coïncide avec la fonctionnelle complète de
Guinand--Weil. L'extinction du terminal et la compatibilité avec les
inclusions permettent de passer à la limite sans introduire de liste de zéros ni le
signe recherché. Par le critère réversible de Weil, tout zéro non trivial
\(\rho\) de la fonction zêta de Riemann satisfait
\[
  \operatorname{Re}\rho=\frac12.
\]
Les calculs régulés à fenêtre finie sont des contrôles ultérieurs ; ils ne font pas
partie des prémisses de la clôture cofinale.

\paragraph{Factorisation cofinale de la conclusion spectrale.}
La chaîne probatoire réunit, dans l'ordre causal déjà établi, la forme
généalogique finie, son complément positif, la naturalité sous les inclusions et
l'extinction compatible du terme terminal. Chaque fenêtre préserve les horloges
premières, le canal archimédien, l'orientation des deux feuillets et le registre de
frontière. La cofinalité transporte conjointement ces données à la limite et
l'identité de Guinand--Weil identifie la forme résultante à sa publication
spectrale classique.

Par conséquent, la positivité ne provient pas d'une sélection ultérieure de
zéros : elle découle de l'opérateur positif construit avant l'évaluation.
L'involution fonctionnelle conserve le feuillet miroir et fixe la droite
\(\operatorname{Re}\rho=1/2\) ; la compatibilité cofinale conserve cette fixation
dans le passage des fenêtres régulées à la forme complète. Les calculs
finis remplissent ainsi leur fonction propre de vérification ultérieure, tandis que
la clôture mathématique réside dans la composition positive et naturelle antérieure.

